% ---------------------------------------------------------------------------
% Results, part B: the bound is attained; the estimator gap; SGN-PINN.
% ---------------------------------------------------------------------------
\subsection{The bound is attained}
\label{sec:res-mc}

A Cram\'er--Rao bound is a statement about a hypothetical efficient estimator
in the asymptotic limit; it becomes a yardstick only if some estimator actually
attains it on this problem. We tested this directly. For each regime, 24
independent datasets (data seeds 101--124, the same design and noise level)
were fitted by the exact maximum-likelihood estimator with $\dP$ known, started
at the truth so that the question is statistical rather than one of finding the
basin (Methods). In the two low-WNT regimes the bound is attained: the
empirical standard deviation of every log-parameter across datasets matches
its CRLB, from $2\times10^{-4}$ to $10^{-2}$ (Fig~\ref{fig:mc}A), the standardised
errors are standard normal (variance 1.03 and 1.13; Fig~\ref{fig:mc}B), the
nominal 95\% intervals $\pm1.96\,s_j$ cover the truth in 95\% and 92\% of cases
(Fig~\ref{fig:mc}C), and twice the drop in cost from the truth to the estimate
follows $\chi^2_{35}$ (mean 36.3 and 33.1; Fig~\ref{fig:mc}D). In the two high-WNT
regimes the estimates scatter less than the bound (variance 0.88 and 0.70,
coverage 95\% and 97\%) and the likelihood-ratio statistic falls short of
$\chi^2_{35}$ (mean 31.9 and 31.0). This is not a failure to converge: restarting
a fit from its own estimate with stricter stopping rules moves it by less than
$10^{-10}$, and tightening the solver moved the likelihood-ratio statistic by
only one to two units. It is the non-asymptotic regime of a weakly informed
likelihood: along the sloppiest directions the log-likelihood is far from
quadratic over the scale of the noise, the estimator is biased there, and its
scatter can fall below a bound that holds for unbiased estimators (the weakest
parameter, $\thP$ in Severe APC Loss, scatters by 33\% against a bound of
57\%). The bound is therefore
attained where the data are informative and conservative where they are not;
in neither case does a converged estimator do worse than it, which is what a
yardstick for estimator performance requires.

\begin{figure}[tbp]
\centering
\fig{fig04_attained.pdf}
\caption{\textbf{The Cram\'er--Rao bound is attained where the data are informative.} Exact maximum-likelihood
estimates ($\dP$ known, started at the truth) on 24 independent datasets per
regime. (A) Empirical standard deviation of each log-parameter across datasets
against its CRLB; grey band, the 95\% sampling interval of a standard deviation
estimated from 24 draws. (B) Standardised errors $(\hat x_j-x_j)/s_j$ pooled
over parameters and datasets, against the standard normal density (dashed). (C)
Coverage of the nominal 95\% interval $\hat x_j\pm1.96\,s_j$ per parameter
(points) and its mean (bar). (D) Twice the cost decrease from the truth to the
estimate against $\chi^2_{35}$ quantiles. All fits use LSODA at relative
tolerance $10^{-11}$ (Methods).}
\label{fig:mc}
\end{figure}

\subsection{Physics-informed estimators leave an estimator gap}
\label{sec:res-gap}

With the yardstick established, the inverse PINN can be placed against it.
Fig~\ref{fig:gap}A compares, on the benchmark data and on an independent
replicate dataset, three PINN variants (each trained once from the common
starting guess) with the classical fit and the bound: the standard formulation
with an automatic-differentiation residual, the same with relative state
weights and a fixed collocation grid, and the same with the derivative-free
integral residual of Eq~\eqref{eq:integral} (Methods). On the benchmark data the
autodiff PINN recovers 14, 10, 7 and 2 of the 34 regular parameters, with
median errors 43 to 92 times the bound. Relative state weights and a fixed
collocation grid lift this to 21, 19, 10 and 10 (10 to 22 times the bound).
Replacing the residual by the integral form on top of them gives 17, 15, 10
and 7 (13 to 33 times), so on these data, at one start each, the
derivative-free residual adds nothing to the weighting. No formulation comes
near the bound (Fig~\ref{fig:gap}A,C).
The integral-residual PINN, the starting point of the SGN-PINN below,
recovers 17, 15, 10 and 7 of the 34 regular
parameters in the four regimes on the benchmark data (18, 16, 13 and 9 on the
replicate), with median relative errors of 8.5\%, 12.2\%, 15.6\% and 30.6\%. On
the same data the converged classical fit recovers all 34 in every regime, and
the CRLB-expected count is 34.0, 34.0, 33.4 and 31.1. Its errors sit at or below
the bound: the median ratio of error to CRLB standard deviation is 1.11, 0.60,
0.21 and 0.42 (an efficient estimator would show 0.67, the median of
$|\mathcal{N}(0,1)|$; values below it in the high-WNT regimes are the
non-asymptotic shrinkage of Section~\ref{sec:res-mc}), whereas the integral PINN's errors
are 13 to 35 times the bound at the median over both datasets (Fig~\ref{fig:gap}B,C). The
shortfall is not concentrated in the poorly informed parameters: parameters
that the design pins to $10^{-3}$ are recovered by the PINN to $10^{-1}$. This
is an estimator gap in the plainest sense: the information is in the data, a
classical estimator extracts it, and the PINN does not.

The classical fit needs care too. From the common starting guess, the
relative-weighted fit reaches the global minimum in every regime, but in Early
Adenoma only after 536 objective evaluations, beyond the cap of 400 used
elsewhere; the unit-weight fit from the same guess converges to a genuine local minimum
(restarting from it with stricter stopping moves it by $10^{-5}$) at a cost 18
times that of the global one (19 of 34 recovered). Every jittered start
reaches the global minimum, and selecting among four starts by cost, which
needs no knowledge of the truth, recovers the bound's count on both datasets
(on the replicate 34, 34, 33 and 31, or 132 against a CRLB-expected 132.5). A few starts are cheap insurance against a likelihood with more
than one basin.

\begin{figure}[tbp]
\centering
\fig{fig05_estimator_gap.pdf}
\caption{\textbf{The estimator gap.} (A) Parameters recovered within 10\% (of the
34 outside the degenerate $(\thP,\dP)$ pair) per regime, for PINNs with the
automatic-differentiation residual, with relative weights and fixed
collocation, and with the integral residual (one start each), the SGN-PINN
(Section~\ref{sec:res-sgn}), the converged classical fit from the common start
and the best of four starts; bars, benchmark data (seed 42); diamonds,
replicate dataset (seed 7; the autodiff and weighted PINN variants were run on
the benchmark data only); hatched, the CRLB-expected count. (B) Absolute
log-error of every regular parameter against its CRLB standard deviation for
the integral PINN and the classical fit, both datasets; solid line, error equal
to the bound; dotted, 10 and 100 times the bound; dashed, 10\% error. (C) Median
ratio of error to bound along the estimator ladder, per regime (both datasets
where available); dashed, an
efficient estimator.}
\label{fig:gap}
\end{figure}

\subsection{Closing the estimator gap: a Schur-complement Gauss--Newton PINN}
\label{sec:res-sgn}

Where does the PINN lose the information? Not in the representation, which
reaches a few percent (Section~\ref{sec:res-forward}), and not in the objective,
whose large-$\lambda$ limit is the classical one; in the optimisation. A
first-order method on a loss whose curvature spans eight orders of magnitude
(Fig~\ref{fig:info}A) moves the well-informed directions quickly and the rest
barely, and the network absorbs what the parameters do not explain. The
SGN-PINN (Methods, Eqs~\eqref{eq:head}--\eqref{eq:schur}) freezes the trained
PINN's hidden features, adds fixed random Fourier features and decaying
exponentials, and optimises only the linear output layer of each condition
together with $\theta$. The inverse problem then becomes a separable nonlinear
least-squares problem whose Gauss--Newton matrix is block-arrowhead
(Fig~\ref{fig:sgn}A), so a full Levenberg--Marquardt step in all
$5.6\times10^4$ unknowns costs ten block factorisations and a $36\times36$
solve. The physics weight $\lambda$ is continued upward from $10^{-2}$, where the
heads mostly fit the data, to $10^3$, where the trajectories must solve the
ODE. Two schedules are run from the integral PINN's estimate, and the reported
one is chosen by the shooting misfit of the ODE solution at its parameters,
which needs no knowledge of the truth. The feature basis and the schedules
were chosen on the benchmark dataset; the replicate dataset is held out.

On the held-out dataset the SGN-PINN recovers 34, 34, 33 and 30 of the 34
regular parameters, with median errors of 0.11\%, 0.20\%, 1.1\% and 1.5\%
(Fig~\ref{fig:gap}A). In Normal, Early and Advanced the continuation converges:
the count rises within the first few iterations at small $\lambda$ and holds as
the physics residual is driven down by three orders of magnitude
(Fig~\ref{fig:sgn}B), and parameter by parameter its errors coincide with those
of the classical fit on the same data (rank correlation of absolute log-errors 0.99--1.00;
Fig~\ref{fig:sgn}C). In Severe APC Loss the physics residual stalls at
$1.3\times10^4$, the trajectory is not an ODE solution, and although 30 of 34
parameters are within 10\% their errors no longer track the classical fit's.
On the benchmark dataset the method converges in Normal and Early (34 and 34)
and stalls in the two high-WNT regimes (28 and 24; residuals $7\times10^3$ and
$1.5\times10^4$). The SGN-PINN therefore reaches the classical estimate from
inside the PINN framework when its continuation converges, and the physics
residual, which needs no truth, says whether it has. In the cases where one
schedule converged to a trajectory far from an ODE solution (physics residual
above $10^5$), the shooting misfit was 0.5--19\% higher and the truth-free
selection discarded it.

The SGN-PINN also returns, as a by-product of its last step, the inverse of the
Schur complement. Its square root tracks the CRLB of the design over three
orders of magnitude, at 1.2 to 1.6 times the bound at the median across regimes
and datasets (Fig~\ref{fig:sgn}D); the excess is expected, because the
relative-weighted objective is itself less efficient than the bound by 7--47\%
(Methods) and the Gauss--Newton curvature at finite $\lambda$ is not exactly the
Fisher information. In an unconverged fit the estimate can be badly wrong (one
parameter in the benchmark Severe fit is reported 2500 times too precise), so
the self-reported precision is only as good as the convergence check that
accompanies it. The method costs 15--27 minutes per regime (both schedules) on one GPU on top
of the PINN it starts from, against 40--43 minutes for that PINN, 2--26 minutes
for one start of the classical fit and under four minutes for the bound itself
(Fig~\ref{fig:sgn}E; all measured on a shared node under load). Its contribution is not speed: it shows that the
estimator gap of a PINN can be closed without abandoning the network
representation of the trajectory, which is what one needs when part of the
dynamics is itself learned, as in the hybrid models below.

\begin{figure}[tbp]
\centering
\fig{fig06_sgn.pdf}
\caption{\textbf{The Schur-complement Gauss--Newton PINN.} (A) Structure of its
Gauss--Newton matrix: dense blocks for the network heads of the ten conditions,
coupled only through the 36 parameters, so each step reduces to the $36\times36$
Schur complement, which is also the profiled curvature that yields standard
errors. (B) Physics residual (top) and parameters within 10\% of 36 (bottom)
along the continuation in $\lambda$ (vertical lines) on the held-out dataset,
per regime, for the schedule chosen by shooting misfit. (C) Absolute log-error
of every regular parameter, SGN-PINN against the classical fit on the same data.
(D) The SGN-PINN's own standard errors against the CRLB of the design. (E)
Measured wall-clock time per regime (range over regimes and datasets; the node
was shared, so times are indicative).}
\label{fig:sgn}
\end{figure}
